% !TeX program = xelatex
% Compile with XeLaTeX in Overleaf.

\documentclass[a4paper]{article}
\usepackage{xeCJK}
\setCJKmainfont[BoldFont=FandolSong-Bold.otf]{FandolSong-Regular.otf}
\usepackage{needspace}

% \usepackage[default]{fontsetup}

\usepackage{fancyhdr}
\usepackage{extramarks}
\usepackage{amsmath}
\usepackage{amsthm}
\usepackage{amsfonts}
\usepackage{tikz}
\usepackage[plain]{algorithm}
\usepackage{algpseudocode}
\usepackage{enumerate}
\usepackage{tikz}

\usetikzlibrary{automata,positioning}

%
% Basic Document Settings
%

\topmargin=-0.2in
\evensidemargin=0in
\oddsidemargin=0in
\textwidth=6.5in
\textheight=9.5in
\headsep=0.25in

\linespread{1.1}

\pagestyle{fancy}
\lhead{}
\chead{\hmwkClass : \hmwkTitle}
\rhead{\firstxmark}
\lfoot{\lastxmark}
\cfoot{\thepage}

\renewcommand\headrulewidth{0.4pt}
\renewcommand\footrulewidth{0.4pt}

\setlength\parindent{0pt}

%
% Create Problem Sections
%

\newcommand{\exitProblemHeader}[1]{
    \nobreak\extramarks{Problem \arabic{#1} }{Problem \arabic{#1} continued on next page\ldots}\nobreak{}
    \stepcounter{#1}
    \nobreak\extramarks{Problem \arabic{#1}}{}\nobreak{}
}

\newcommand*\circled[1]{\tikz[baseline=(char.base)]{
    \node[shape=circle,draw,inner sep=2pt] (char) {#1};}}

\setcounter{secnumdepth}{0}
\newcounter{partCounter}
\newcounter{homeworkProblemCounter}
\setcounter{homeworkProblemCounter}{1}
\nobreak\extramarks{Problem \arabic{homeworkProblemCounter}}{}\nobreak{}

%
% Homework Problem Environment
%

\newenvironment{homeworkProblem}[1][-1]{
    \ifnum#1>0
        \setcounter{homeworkProblemCounter}{#1}
    \fi
    \section{Problem \arabic{homeworkProblemCounter}}
    \setcounter{partCounter}{1}
}{
    \exitProblemHeader{homeworkProblemCounter}
}

%
% Homework Details
%

\newcommand{\hmwkTitle}{Homework\ \#04}
\newcommand{\hmwkClass}{Probability \& Statistics for EECS}
\newcommand{\hmwkDueDate}{2026-10-27}
\newcommand{\hmwkAuthorName}{杨涵}
\newcommand{\hmwkAuthorID}{2025531064}

%
% Title Page
%

\title{
    \vspace{2in}
    \textmd{\textbf{\hmwkClass:\\ \hmwkTitle}}\\
    \normalsize\vspace{0.1in}\small{Due\ on\ \hmwkDueDate\ at 18:00}\\
    \vspace{4in}
}

\author{
    Name: \textbf{\hmwkAuthorName} \\
    Student ID: \hmwkAuthorID}
\date{}

\renewcommand{\part}[1]{\textbf{\large Part \Alph{partCounter}}\stepcounter{partCounter}\\}

%
% Various Helper Commands
%

\newcommand{\alg}[1]{\textsc{\bfseries \footnotesize #1}}
\newcommand{\deriv}[1]{\frac{\mathrm{d}}{\mathrm{d}x} (#1)}
\newcommand{\pderiv}[2]{\frac{\partial}{\partial #1} (#2)}
\newcommand{\dx}{\mathrm{d}x}
\newcommand{\solution}{\par\addvspace{\baselineskip}\Needspace{4\baselineskip}\noindent\textbf{\large Solution}\par\nobreak\smallskip}
\newcommand{\E}{\mathrm{E}}
\newcommand{\Var}{\mathrm{Var}}
\newcommand{\Cov}{\mathrm{Cov}}
\newcommand{\Bias}{\mathrm{Bias}}

\begin{document}

\date{
    Due on Oct. 27, 2026, 18:00 UTC+8}
\title{SI 140A-01  Probability \& Statistics for EECS, Fall 2026 \\
    Homework 4}
\maketitle
Read all the instructions below carefully before you start working on the assignment, and before you make a submission.
\begin{itemize}
    \item You are required to write down all the major steps towards making your conclusions; otherwise you may obtain limited points of the problem.
    \item Write your homework in English; otherwise you will get no points of this homework.
    \item Any form of plagiarism will lead to $0$ point of this homework.
\end{itemize}

\newpage

\begin{homeworkProblem}
A book has $n$ typos. Two proofreaders, Prue and Frida, independently read the book.
Prue catches each typo with probability $p_1$ and misses it with probability $q_1 = 1-p_1$, independently, and likewise for Frida, who has probabilities $p_2$ of catching and $q_2=1-p_2$ of missing each typo. Let $X_1$ be the number of typos caught by Prue, $X_2$ be the number caught by Frida, and $X$ be the number caught by at least one of the two proofreaders.
\begin{enumerate}
    \item[(a)] Find the distribution of $X$.
    \item[(b)] For this part only, assume that $p_1 = p_2$. Find the conditional distribution of $X_1$ given that $X_1 + X_2 = t$.
\end{enumerate}

\solution
\textbf{(a)} The probability that both proofreaders miss a typo is
\[
q_1q_2=(1-p_1)(1-p_2).
\]
Thus,
\[
X\sim\mathrm{Bin}(n,1-q_1q_2),
\]
with
\[
P(X=k)=\binom{n}{k}(1-q_1q_2)^k(q_1q_2)^{n-k},
\qquad k=0,1,\dots,n.
\]

\medskip
\Needspace{17\baselineskip}
\textbf{(b)} Let $p_1=p_2=p$ and write $q=1-p$. Since $X_1$ and $X_2$
are independent,
\[
X_1\sim\mathrm{Bin}(n,p),\qquad
X_2\sim\mathrm{Bin}(n,p),\qquad
X_1+X_2\sim\mathrm{Bin}(2n,p).
\]
Given $X_1+X_2=t$, we have
\begin{align*}
P(X_1=k\mid X_1+X_2=t)
&=\frac{P(X_1=k,\ X_2=t-k)}{P(X_1+X_2=t)}\\
&=\frac{P(X_1=k)P(X_2=t-k)}{P(X_1+X_2=t)}\\
&=\frac{\binom{n}{k}(1-q)^kq^{n-k}
        \binom{n}{t-k}(1-q)^{t-k}q^{n-t+k}}
       {\binom{2n}{t}(1-q)^tq^{2n-t}}\\
&=\frac{\binom{n}{k}\binom{n}{t-k}}{\binom{2n}{t}}.
\end{align*}
Therefore,
\[
P(X_1=k\mid X_1+X_2=t)
=\frac{\binom{n}{k}\binom{n}{t-k}}{\binom{2n}{t}},
\qquad 0\leq k\leq\min(n,t).
\]

\end{homeworkProblem}

\newpage

\begin{homeworkProblem}
Let $X$ be a random day of the week, coded so that Monday is $1$, Tuesday is $2$, \textit{etc}. (so $X$ takes values $1, 2, \dots, 7$ with equal probabilities).
Let $Y$ be the next day after $X$.
Do $X$ and $Y$ have the same distribution?
What is $P(X < Y)$?

\solution
We have
\[
P(X=k)=\frac17,\qquad P(Y=k)=\frac17,
\qquad k=1,2,\dots,7.
\]
Thus, $X$ and $Y$ have the same distribution. The possible pairs are
\[
\begin{array}{c|ccccccc}
X & 1 & 2 & 3 & 4 & 5 & 6 & 7\\
\hline
Y & 2 & 3 & 4 & 5 & 6 & 7 & 1
\end{array}
\]
The inequality $X<Y$ holds for the first six pairs, so
\[
P(X<Y)=\frac67.
\]

\end{homeworkProblem}

\newpage

\begin{homeworkProblem}
For $X$ and $Y$ binary digits (0 or 1), let $X \oplus Y$ be 0 if $X=Y$ and 1 if $X\neq Y$ (this operation is called exclusive or (often abbreviated to XOR), or addition mod 2).
\begin{enumerate}[(a)]
    \item Let $X \sim {\rm Bern}(p)$ and $Y \sim {\rm Bern}(1/2)$, independently. What is the distribution of $X \oplus Y$?
    \item With notation as in sub-problem (a), is $X \oplus Y$ independent of $X$?  Is $X \oplus Y$ independent of $Y$?  Be sure to consider both the case $p= 1/2$ and the case $p\neq 1/2$.
    \item Let $X_1, ... ,X_n$ be i.i.d. (i.e., independent and identically distributed) Bern(1/2) R.V.s. For each nonempty subset $J$ of $\{1,2,...,n\}$, let
    \begin{equation*}
        Y_J=\oplus_{j \in J}X_j.
    \end{equation*}
    Show that $Y_J \sim {\rm Bern}(1/2)$ and that these $2^n-1$ R.V.s are pairwise independent, but not independent.
\end{enumerate}

\solution
\textbf{(a)} Let $Z=X\oplus Y$. The truth table is
\[
\begin{array}{cc|c}
X & Y & Z\\
\hline
0 & 0 & 0\\
0 & 1 & 1\\
1 & 0 & 1\\
1 & 1 & 0
\end{array}
\]
Since $X\sim\mathrm{Bern}(p)$ and $Y\sim\mathrm{Bern}(1/2)$ independently,
\[
P(X=1)=p,\qquad P(X=0)=1-p,\qquad
P(Y=1)=P(Y=0)=\frac12.
\]
Hence,
\begin{align*}
P(Z=1)
&=P(X=0,Y=1)+P(X=1,Y=0)
=\frac12(1-p)+\frac12p=\frac12,\\
P(Z=0)
&=P(X=1,Y=1)+P(X=0,Y=0)
=\frac12p+\frac12(1-p)=\frac12.
\end{align*}
Thus, $X\oplus Y\sim\mathrm{Bern}(1/2)$.

\medskip
\Needspace{15\baselineskip}
\textbf{(b)} If $X$ is known, then $Z$ is either $Y$ or $1-Y$. In either
case, the probabilities of $Z=0$ and $Z=1$ remain $1/2$. Therefore,
$X\oplus Y$ is independent of $X$.

\medskip
If $Y$ is known, then
\[
P(Z=1\mid Y=0)=P(X=1)=p,\qquad
P(Z=1\mid Y=1)=1-p.
\]
\begin{enumerate}[(i)]
    \item If $p=1/2$, then $X\oplus Y$ is independent of $Y$.
    \item If $p\neq1/2$, then $X\oplus Y$ is not independent of $Y$.
\end{enumerate}

\newpage

\textbf{(c)} For a nonempty subset $J$, we have
\[
Y_J=\bigoplus_{j\in J}X_j.
\]
Fix all the variables in $Y_J$ except one, $X_j$. Then $Y_J$ is either
$X_j$ or $1-X_j$. Since $X_j\sim\mathrm{Bern}(1/2)$, it follows that
$Y_J\sim\mathrm{Bern}(1/2)$.

\medskip
\Needspace{8\baselineskip}
Let $J$ and $K$ be two distinct nonempty subsets. After interchanging their
names if necessary, choose an index $j\in J\setminus K$. Fix all the
variables except $X_j$. Then $Y_K$ is fixed, while $Y_J$ is either $X_j$
or $1-X_j$. Since $X_j\sim\mathrm{Bern}(1/2)$, $Y_J$ still has distribution
$\mathrm{Bern}(1/2)$. Thus, knowing one of $Y_K$ and $Y_J$ does not change
the distribution of the other, so they are pairwise independent.

\medskip
\Needspace{4\baselineskip}
However, they are not mutually independent. Consider $X_1$, $X_2$, and
$X_3=X_1\oplus X_2$; these variables are dependent.

\end{homeworkProblem}

\newpage

\begin{homeworkProblem}
Let $X$ be a discrete R.V. whose distinct possible values are $a_1, a_2,...,a_n,$ with probabilities $p_1, p_2 ...,p_n,$ respectively (so $p_1 + p_2 + \cdots + p_n = 1$). The entropy of $X$ is defined to be the average surprise of learning the value of $X$:
\begin{equation*}
    H(X)=\sum_{j=1}^np_j\log_2(1/p_j).
\end{equation*}
Show that the maximum entropy for $X$ is when
its distribution is uniform over $a_1, a_2,...,a_n,$ i.e., $p_j = 1/n, \forall j$.

\solution
Assume $p_j>0$. For the concave function $f(x)=\log_2 x$, Jensen's
inequality gives
\[
\sum_{i=1}^{n}a_i f(x_i)
\leq f\left(\sum_{i=1}^{n}a_i x_i\right).
\]
Therefore,
\[
H(X)=\sum_{j=1}^{n}p_j\log_2\left(\frac1{p_j}\right)
\leq\log_2\left(\sum_{j=1}^{n}p_j\frac1{p_j}\right)
=\log_2 n.
\]
Equality holds if and only if
\[
p_j=\frac1n,\qquad \forall j.
\]

\end{homeworkProblem}

\end{document}
